\par\Needspace{6\baselineskip}\medskip
\begin{center}\hypertarget{native-180-article134}{}\pdfbookmark[2]{134.}{article-134}134.\end{center}
From the preceding discussion it will now be clear why in our method we take for $x,y$,
respectively, the quantities
\[
\frac{n''}{n}=P\qquad\text{and}\qquad
2\left(\frac{n+n''}{\mathpunct{\mkern-12mu.\mkern12mu}n'}-1\right)r'^3=Q
\]
It is clear, \textit{first}, that if $P$ and $Q$ are regarded as known, $\delta'$ can be
determined from them by the equation
\[
a\delta'=b+\frac{c+dP}{1+P}\left(1+\frac{Q}{2r'^3}\right)
\]
and then $\delta$ and $\delta''$ by equations 4 and 6 of Article 114, since we have
\[
\frac{n}{n'}=\frac{1}{1+P}\left(1+\frac{Q}{2r'^3}\right)\mathpunct{\raisebox{.45ex}{,}}\qquad
\frac{n''}{n'}=\frac{P}{1+P}\left(1+\frac{Q}{2r'^3}\right)\mathpunct{\cdot}
\]
\textit{Second}, it is clear that in the first hypothesis, for the quantities $P,Q$ whose exact
values are
\hypertarget{native-181}{}
\[
\frac{\theta''}{\theta}\cdot\frac{\eta}{\eta''}\mathpunct{\raisebox{.45ex}{,}}\qquad
\frac{r'r'\theta\theta''}{rr''\eta\eta''\cos f\cos f'\cos f''}\mathpunct{\raisebox{.45ex}{,}}
\]
the approximate values $\dfrac{\theta''}{\theta}\mathpunct{\raisebox{.45ex}{,}}\;\theta\theta''$ are available at once. This hypothesis will introduce errors of
the first order into the determination of $\delta'$ and hence also of $\delta,\delta''$, or errors of the
second order in the special case mentioned several times. Although these conclusions can
generally be relied on with complete confidence, they may lose their force in one particular
case. This case arises whenever the quantity $(0.1.2)$, which is generally of the third order,
happens to become $=0$ or is so small that it must be assigned to a higher order. This occurs
whenever the geocentric motion on the celestial sphere exhibits a point of inflection near the
middle position. Finally, it is clear that our method requires the heliocentric motion over the
interval covered by the three observations to be sufficiently small. The very complicated nature
of the problem makes this restriction unavoidable. Nor need it be considered a disadvantage,
since one will always wish to undertake the first determination of the unknown orbit of a new
celestial body as soon as possible. Moreover, the restriction can be understood fairly broadly,
as the examples given below will show.
